En este informe se presenta el montaje y validación experimental de un conversor digital a analógico de 3 bits. El marco teórico del circuito se fundamentó en la aplicación del teorema del equivalente de Thévenin para determinar la contribución de voltaje individual de cada bit y en el principio de superposición para calcular la tensión de salida total ante combinaciones binarias de entrada. El montaje práctico se realizó en protoboard empleando 11 resistencias de $(1\pm0.05)$ \unit{\kilo\ohm} y un voltaje de de fuente de 16 \unit{\volt} DC. Las mediciones registradas para las ocho combinaciones posibles (000 a 111) mostraron un comportamiento lineal en el rango de 0 a 7 \unit{\volt}, presentando una coincidencia notable con el modelo ideal y un error relativo porcentual máximo de apenas 40\%. Las pequeñas desviaciones experimentales se justifican por la tolerancia nominal de los componentes, pequeñas inestabilidades al ajustar la fuente DC y la resolución de escala del multímetro digital.
